\section{Discussão dos resultados}

Esta discussão utiliza exclusivamente os dados registrados em
\texttt{resultados\_completos.csv}, as tabelas e gráficos gerais gerados a partir
desse arquivo, o dashboard estático e os dados detalhados dos gráficos de Gantt.
Assim, as interpretações apresentadas estão restritas ao conjunto de instâncias
executado, às formulações implementadas e à configuração computacional adotada:
limite de tempo de 300 segundos por instância e tolerância relativa de gap de
1\%.

\subsection{Resultados observados}

Foram analisadas 48 instâncias no total: 31 de Machine Scheduling, 10 de Job
Shop Scheduling e 7 de Flow Shop Scheduling. O HiGHS encontrou solução primal
viável em todas as instâncias registradas, e não houve erros de execução nos
resultados consolidados.

O problema com maior proporção de soluções ótimas foi o Flow Shop Scheduling,
com 6 soluções ótimas em 7 instâncias, correspondendo a aproximadamente 85,7\%
do conjunto desse problema. O Job Shop Scheduling apresentou 6 soluções ótimas
em 10 instâncias, isto é, 60,0\%. O Machine Scheduling teve 17 soluções ótimas
em 31 instâncias, aproximadamente 54,8\%. Portanto, embora Machine Scheduling
tenha apresentado a maior quantidade absoluta de soluções ótimas, Flow Shop
Scheduling apresentou a maior proporção de instâncias resolvidas dentro da
tolerância de otimalidade configurada.

Em relação ao tempo de solução, o Job Shop Scheduling apresentou o maior tempo
médio, 190,464 segundos, e também o maior tempo mediano, 235,695 segundos. O
Machine Scheduling teve tempo médio de 150,670 segundos e mediana de 119,012
segundos. O Flow Shop Scheduling apresentou o menor tempo médio, 44,178
segundos, e a menor mediana, 0,619 segundo. Esses valores indicam que, no
conjunto de dados considerado, Job Shop foi o problema com comportamento mais
custoso em termos de tempo central típico, enquanto Machine Scheduling concentrou
maior número absoluto de interrupções por limite de tempo.

As instâncias que atingiram o limite de tempo foram 14 em Machine Scheduling, 4
em Job Shop Scheduling e 1 em Flow Shop Scheduling. Em Machine Scheduling, todas
as interrupções ocorreram em instâncias com 15 ou mais tarefas. Em Job Shop, as
instâncias interrompidas foram \texttt{abz5}, \texttt{ft10}, \texttt{ft20} e
\texttt{orb01}. Em Flow Shop, apenas \texttt{problem\_10m\_20j} atingiu o limite
de tempo. O dashboard e a tabela de instâncias com maior tempo mostram que
essas instâncias aparecem entre os casos mais custosos do experimento.

O comportamento do gap também variou substancialmente entre os problemas. O gap
médio geral foi 0,3841 em Machine Scheduling, 0,1297 em Job Shop Scheduling e
0,0079 em Flow Shop Scheduling. Considerando apenas as soluções não ótimas, os
gaps médios foram 0,8432 para Machine Scheduling, 0,3118 para Job Shop
Scheduling e 0,0241 para Flow Shop Scheduling. O maior gap registrado ocorreu em
Machine Scheduling, com valor aproximadamente igual a 1,0000. Em Job Shop, o
maior gap foi 0,6008, observado em \texttt{ft20}. Em Flow Shop, o maior gap foi
0,0241, correspondente à única instância interrompida pelo limite de tempo.

\subsection{Possíveis explicações}

Os resultados sugerem que a dificuldade computacional está associada à
combinação entre tamanho da instância, número de variáveis, número de restrições
e estrutura combinatória do modelo. Essa relação aparece nos gráficos de
tarefas versus tempo, variáveis versus tempo e restrições versus tempo, embora
não seja perfeitamente determinística.

No Machine Scheduling, todas as instâncias possuem apenas uma máquina, mas o
modelo utiliza variáveis binárias associadas a pares de tarefas. A média de
variáveis foi 148,548 e a média de restrições foi 445,645, com máximos de 348
variáveis e 1044 restrições. As instâncias de maior porte desse problema,
especialmente as com 21 e 24 tarefas, aparecem entre as maiores em número de
variáveis, restrições e gap. Uma explicação plausível é que o crescimento
combinatório das decisões de sequenciamento dificulta o fechamento do gap, mesmo
em máquina única.

No Job Shop Scheduling, o número médio de tarefas foi 10,6 e o número médio de
máquinas foi 7,1. A média de variáveis foi 448,6 e a média de restrições foi
1344,8, superiores às médias de Machine Scheduling. O maior modelo de Job Shop
registrado teve 1051 variáveis e 3152 restrições. Como o Job Shop combina
precedências tecnológicas por job com conflitos de capacidade em várias
máquinas, é razoável interpretar que a estrutura disjuntiva do modelo contribuiu
para os tempos elevados. Essa interpretação é compatível com o fato de o Job
Shop ter apresentado o maior tempo médio e mediano.

No Flow Shop Scheduling, a formulação posicional explorou a existência de uma
sequência comum de jobs nas máquinas. Apesar de o conjunto incluir instâncias
com até 30 tarefas e 10 máquinas, o problema apresentou a menor média de tempo e
a maior proporção de soluções ótimas. A média de variáveis foi 439,571 e a média
de restrições foi 741,286. Assim, embora os modelos de Flow Shop não sejam
pequenos em número de variáveis, a estrutura regular da sequência comum pode ter
favorecido o desempenho observado do HiGHS neste conjunto específico.

O número de máquinas também afeta os problemas de maneiras distintas. Em Machine
Scheduling ele é constante e igual a 1, logo a dificuldade observada está ligada
principalmente ao número de tarefas e às disjunções entre elas. Em Job Shop, o
número de máquinas aumenta a quantidade de conflitos de capacidade entre
operações. Em Flow Shop, mais máquinas aumentam a dimensão do modelo e a cadeia
de precedências, mas preservam a ordem comum das tarefas, o que pode explicar
por que instâncias de Flow Shop com várias máquinas ainda foram resolvidas com
mais estabilidade do que parte das instâncias de Job Shop.

\subsection{Análise dos Gantts}

Os gráficos de Gantt foram gerados para três instâncias representativas:
\texttt{inst\_n05\_s01} em Machine Scheduling, \texttt{ft06} em Job Shop
Scheduling e \texttt{problem\_3m\_10j} em Flow Shop Scheduling. Essas três
soluções foram resolvidas novamente apenas para recuperar os horários detalhados
de início e término, e as validações estruturais não indicaram violações de
precedência, duração, sobreposição ou makespan.

No Gantt de Machine Scheduling, a instância possui 5 tarefas e uma única
máquina. O makespan da programação detalhada foi 39,0. A soma dos tempos de
processamento foi 38,0, e a máquina ficou ociosa por 1,0 unidade de tempo quando
se considera o intervalo entre o tempo zero e o makespan. No intervalo efetivo
entre a primeira tarefa e a última conclusão, não houve ociosidade interna. Como
há apenas uma máquina, ela concentra toda a carga e determina diretamente o
makespan. O gráfico evidencia uma sequência compacta de tarefas após o início
da primeira operação.

No Gantt de Job Shop, a instância \texttt{ft06} possui 6 jobs, 6 máquinas e 36
operações, com makespan aproximadamente igual a 55,0. Todas as operações
esperadas estão presentes. As precedências de cada job foram respeitadas, mas
existem esperas entre operações: foram observados 30 intervalos entre operações
consecutivas de jobs, dos quais 14 tiveram espera positiva, com soma total de
esperas igual a 53,0 e maior espera igual a 8,0. Essa característica é coerente
com o Job Shop, no qual uma operação pode precisar aguardar a liberação da
máquina seguinte. Pelos dados do Gantt, as máquinas M5, M0 e M4 apresentaram
altas cargas de processamento, com 43,0, 40,0 e 40,0 unidades, respectivamente.
A máquina M5 teve a menor folga em relação ao makespan, 12,0 unidades, o que
sugere papel relevante na formação do caminho crítico da solução observada. As
máquinas M2 e M3 apresentaram maiores períodos ociosos dentro de suas janelas de
atividade, indicando que nem todas as máquinas foram igualmente pressionadas.

No Gantt de Flow Shop, a instância \texttt{problem\_3m\_10j} possui 10 tarefas e
3 máquinas, com makespan 126,0. A sequência encontrada foi
\texttt{J4 -> J8 -> J9 -> J3 -> J5 -> J7 -> J10 -> J1 -> J2 -> J6}. O fluxo das
tarefas respeita a ordem M1, M2 e M3. As cargas totais por máquina foram 101,0
em M1, 124,0 em M2 e 101,0 em M3. A máquina M2 foi a mais carregada e terminou
em 125,0, muito próxima do makespan, o que sugere que ela atuou como gargalo
principal nesta programação. A máquina M3 iniciou apenas no tempo 16,0 e teve
ociosidade interna de 9,0 unidades dentro de sua janela de atividade, além de
receber tarefas condicionadas pela conclusão em M2. Foram observadas esperas
entre máquinas em 13 dos 20 intervalos analisados, com soma total de 129,0 e
maior espera igual a 18,0. Esse comportamento é compatível com o fluxo em série:
mesmo com sequência comum, as máquinas finais podem acumular esperas decorrentes
do desalinhamento entre tempos de processamento nas etapas anteriores.

\subsection{Desempenho do HiGHS e limites da configuração}

O HiGHS foi capaz de encontrar soluções viáveis para todas as 48 instâncias
registradas. Esse é um resultado relevante, pois mesmo as instâncias
interrompidas por tempo possuem ponto primal viável. Por outro lado, a
configuração de 300 segundos e gap relativo de 1\% não foi suficiente para
encerrar todas as instâncias dentro da tolerância estabelecida. O efeito dessa
configuração é visível principalmente em Machine Scheduling e Job Shop
Scheduling, que concentraram a maior parte dos limites de tempo e dos gaps mais
altos.

Também é importante observar que o status \texttt{OPTIMAL} deve ser interpretado
dentro da tolerância configurada. Algumas instâncias com status ótimo possuem
gap positivo menor que 1\%, o que indica otimalidade de acordo com o critério de
parada adotado. Portanto, a comparação entre problemas deve considerar não
apenas o status final, mas também o gap, o tempo consumido e o tamanho do modelo.

\subsection{Limitações do experimento}

As conclusões são limitadas ao conjunto de instâncias analisado. Os tamanhos dos
conjuntos não são iguais: foram 31 instâncias de Machine Scheduling, 10 de Job
Shop Scheduling e 7 de Flow Shop Scheduling. Assim, médias e proporções devem
ser interpretadas com cautela, principalmente no Flow Shop, que possui menor
número de observações.

Outra limitação é que as formulações utilizadas são modelos MILP diretos. Em
Machine Scheduling e Job Shop, as formulações dependem de restrições
disjuntivas e variáveis binárias de ordenação, que podem produzir relaxações
lineares difíceis para o solver. No Flow Shop, a formulação posicional pareceu
mais favorável nos dados observados, mas isso não implica superioridade geral do
problema ou da formulação em outros conjuntos de instâncias.

Por fim, os resultados consolidados armazenam métricas agregadas, mas não
armazenam as programações completas de todas as soluções. Por isso, a validação
estrutural detalhada foi realizada apenas para as três instâncias selecionadas
para os gráficos de Gantt. Para análises futuras, seria recomendável registrar
também os horários de início e término de todas as operações, permitindo validar
e visualizar qualquer solução sem necessidade de nova execução do modelo.
